\chapter{Администратору}
\label{ch:admin}

Глава для тех, у кого в профиле стоит признак администратора --- в шапке сайта рядом с именем видна корона.

\section{Группы доступа}
\label{sec:admin:groups}

Раздел \texttt{/access\_groups} решает одну задачу: кто какие сценарии видит. Страница разделена на две вкладки.

\paragraph{Вкладка ``Группы''}
Список групп с поиском по названию. Кнопка ``Добавить'' создаёт новую --- сайт спросит название. Кнопка ``Настройки'' у каждой группы открывает её карточку (\texttt{/access\_groups/<id>}), где задаётся состав участников и доступные сценарии.

\paragraph{Вкладка ``Пользователи''}
Список всех пользователей платформы с поиском по имени, почте и Telegram-идентификатору. Значки рядом с именем показывают роли. По клику открывается профиль пользователя и его группы.

\begin{quote}
\small
Учтите: проверка прав стоит только на списке групп. Прямая ссылка на карточку конкретной группы откроется у любого авторизованного пользователя. Пока это не исправлено, не считайте адреса карточек секретными.
\end{quote}

\section{Выдача прав мастера}
\label{sec:admin:grant-master}

Интерфейса для этого пока нет --- ни на странице пользователя, ни в группах. Признак ``может быть мастером'' проставляется прямо в базе:

\begin{quote}
\ttfamily\small
UPDATE "user" SET can\_be\_master = true WHERE email = 'master@example.com';
\end{quote}

Тем же способом выдаётся признак администратора --- поле \texttt{is\_admin}. После изменения пользователю нужно перезайти на сайт, чтобы новые права подхватились.

\paragraph{Первый администратор}
При разворачивании платформы с нуля администраторов нет вообще, поэтому первого приходится назначать этим же запросом. Дальше он может пользоваться разделом групп доступа.

\section{Обязательное при разворачивании}
\label{sec:admin:deploy}

Одна вещь, без которой платформу нельзя выпускать наружу: переменная окружения \texttt{SECRET\_KEY}. Ею подписываются токены входа. Значение по умолчанию лежит в открытом репозитории --- с ним кто угодно подделает вход под любым пользователем.

В рабочем режиме (\texttt{ENV=production}) приложение не запустится, если ключ не задан или оставлен стандартным. В режиме разработки оно запустится, но напишет предупреждение в лог.
